\section{Conclusion and Future Work}

The key finding is that despite the wide disparity in local accounting tree structures across 23 national jurisdictions, a \textbf{minimum universal core of 18 accounts} (Level 3 nodes) anchors all analyzed statutory systems. Building Kontablo around this core—rather than attempting a politically-negotiated compromise—offers a practical path toward the ultimate goal: a 99\%-automatable global accounting protocol, heavily supervised and ultimately authorized by human expert oversight.

Our analysis demonstrates that the 30 Level 3 accounts cover **92\% of routine business transactions by empirical volume (count)**—as benchmarked against thousands of generic SME ledger exports—validating that extreme ledger granularity in national standards is mostly reserved for rare statistical or bureaucratic edge cases.

Future work for Kontablo includes: 
\begin{enumerate}
    \item \textbf{Expert Validation Phase (Q2 2026):} Conducting a series of structured validation interviews with 15 expert CPAs from diverse OECD and non-OECD nations.
    \item \textbf{Industry-Specific Extensions:} Developing specialized ontology overlays for banking (IFRS 9/7), insurance (IFRS 17), and agriculture (IAS 41).
    \item \textbf{Deployment of Universal ERP Connectors:} Deploying production-grade two-way API connectors for the 10 primary commercial systems evaluated in section 5.6 (SAP, NetSuite, Dynamics 365, etc.).
    \item \textbf{High-Frequency Agentic Tracking:} Standardizing \texttt{Agent\_ID} fields and batching mechanisms within the Kontablo Graph to track liability and authorization trails for the burgeoning autonomous internet commerce economy.
\end{enumerate}

The Kontablo Universal Accounting Ontology is currently published as an open-source project under an MIT license, with the full research data, mass-consolidation engine, and AI mapping services available for public auditing and collaborative extension.
